We pursue \emph{semi-supervised} learning (SSL)~\citep{zhuSemiSupervisedLearningLiterature2005,vanengelenSurveySemisupervisedLearning2020} for the specific problem of image classification with deep neural networks~\citep{oliverRealisticEvaluationDeep2018}.
We can observe images $x$ (represented as $D$-dimensional vectors) as well as corresponding labels $y \in \{1, 2, \ldots C\}$ for $C$ classes of interest.
The goal is to train predictors from both a labeled dataset $\mathcal{D}^L$ of feature-labeled pairs $x,y$ and an unlabeled dataset $\mathcal{D}^U$ containing feature vectors $x$ only.

%% MCH: Cut this bc too redundant with earlier
%In our motivating applications, unlabeled examples are easily collected, while labels are difficult due to burdensome requirements of time, money, or human expertise.
%Medical image classification is well-suited to SSL, because routinely collected images are available in abundance, but annotation costs usually mean only small labeled sets are affordable. 

%%%%%% TABLE: RELATED WORK
\setlength{\tabcolsep}{.24cm}
\begin{table*}[!t]
\begin{tabular}{l r  l l l}
Method & Acc. & Paradigm 
	%& \todo{WHAT ABOUT DELETING THIS COLUMNApplicability} 
	& Extra Complexity & Realistic Eval.
\\ \hline
Fix-A-Step (ours) 
	& 85.4 % 85.45 verified in google doc, MT_CombinedAugGrad
	& OOD \emph{helpful}
	%& any $\ell^U$
	& none
	& Heart2Heart
\\
TOOR {\tiny \citep{huangTheyAreNot2022}}
	& $*$78.5 % read from Fig 4 of Huang et al TOOR paper
	& OOD \emph{helpful}
	%& any $\ell^U$
	& Separate NN for OOD discrimination
	& none
\\
CL {\tiny{\citep{cascante-bonillaCurriculumLabelingRevisiting2021}}}
	& $*$83.0
	& OOD harmful % Quote from paper: "We conjecture that our self-pacing curriculum is key to this scenario, where the adaptive thresholding scheme could help filter the out-of-distribution unlabeled samples during training"
	%& pseudo-label $\ell^U$
	& Multiple rounds of training, each from scratch
	& none
\\
OpenMatch {\tiny{\citep{saitoOpenMatchOpensetConsistency2021}}}
	& 82.3 % 82.25
	& OOD harmful
	%& \todo{FixMatch only?}
	& Extra one-vs-all OOD detector / class
	& none
\\
DS3L {\tiny \citep{guoSafeDeepSemiSupervised2020}}
	& 74.7 %74.7 from our implementation 
	% in google doc, Row "DS3L + MT (reported)"
	% sanity check: looks like 77.5 in their Fig 6
	& OOD harmful
	%& any $\ell^U$
	& 3x train time due to bilevel optimization
	& none
\\
MTCF {\tiny \citep{yuMultiTaskCurriculumFramework2020}}
	& 77.0 % read from our Fig 1 (after adding cosine LR, 68.0-->77.0)
	& OOD harmful
	%& any $\ell^U$
	& Extra OOD head, curriculum learning
	& none
\\
UASD {\tiny \citep{chenSemiSupervisedLearningClass2020}}
	& $*$78.2 % in google doc, row "UASD Figure 3 left (reported)"
	& OOD harmful
	%& fixed $\ell^U$
	& none
	& none
\\
Safe-Student
	{\tiny{\citep{heSafeStudentSafeDeep2022}}}
	& $^{\dagger}$n/a
	& OOD harmful
	%& \todo{FILL} %their objective eq11 doesnt fit our eq1?
	& 2 NNs (teacher \& student), extra KL loss
	& none
\end{tabular}
\caption{Comparison of related work on open-set/safe SSL.
\emph{Acc} means accuracy on the CIFAR-10 6-animal task (defined in Sec.~\ref{sec:results-cifar}) with 400 labeled examples/class and an open-set unlabeled set (50\% mismatch).
Fix-A-Step uses a FixMatch base model, as does OpenMatch.
%, or MeanTeacher for older methods.
Numbers come from our implementation except if marked \textbf{*} (copied from cited paper) or $^\dagger$ (not assessed in cited paper).
%\emph{Time:} Runtime for training on 6-animal task.
\emph{Paradigm:} how each method treats out-of-distribution (OOD) images in the unlabeled set, broadly either possibly helpful or likely harmful (and thus in need of filtering).
%\emph{Applicability:} whether the method might be possible with different unlabeled losses $\ell^U$.
\emph{Extra Complexity:} additional neural networks, layers, or runtime concerns that exceed a standard SSL deep classifier like MixMatch.
\emph{Realistic Eval.:} evaluation beyond ``artificial'' unlabeled sets from common datasets like CIFAR, ImageNet, etc.
%$\dagger$: \todo{DELETE? OpenMatch has no realistic SSL evaluation, but evaluates OOD detection on fine-grained data (e.g. breeds of dog).}
%\todo{CITE} and Dogs\todo{CITE}.
}%endcaption	
\vspace{-0.1cm}
\label{tab:related_work_comparison}
\end{table*}



\paragraph{Training for Deep SSL.}
While many SSL paradigms have been tried~\citep{kingmaSemisupervisedLearningDeep2014,kumarSemisupervisedLearningGANs2017,nalisnickHybridModelsDeep2019}, the dominant approaches for semi-supervised training of deep image classifiers today continue to modify standard objectives for discriminative neural nets by adding a regularization term using unlabeled data~\citep{miyatoVirtualAdversarialTraining2019,sohnFixMatchSimplifyingSemisupervised2020}.
This approach trains a neural net probabilistic classifier $f$ with weights $w$ by solving the optimization problem:
\begin{align} % Use \! to reduce whitespace between symbols
\min_{w} ~ \!
	\textstyle \sum_{x,y \in \mathcal{D}^L} \!
		\ell^L( y, f_w(x) )	
	+ 
	\lambda \!
	\sum_{x \in \mathcal{D}^U} \!
		\ell^U(x; w)
	\label{eq:loss-for-consistency-ssl}
\end{align}
Here, the first term $\ell^L$ is a labeled-set-only \emph{cross entropy} loss and the second term $\ell^U$ is a method-specific unlabeled-set loss.
A key hyperparameter is the unlabeled-loss-weight $\lambda > 0$, which balances the two terms.
Approaches such as the Pi-model~\citep{laineTemporalEnsemblingSemiSupervised2017}, Pseudo-Label~\citep{leePseudoLabelSimpleEfficient2013}, Mean-Teacher~\citep{tarvainenMeanTeachersAre2017}, Virtual Adversarial training (VAT)~\citep{miyatoVirtualAdversarialTraining2019}, and FixMatch~\citep{sohnFixMatchSimplifyingSemisupervised2020} all fit this objective, with variations in (1) the choice of function for $\ell^U$; (2) how data augmentation may alter images $x$; (3) procedures within $\ell^U$ that produce a perturbed image $x'$ that should be consistent with $x$; and (4) optimization routines to solve for $w$.

\textbf{Uncurated SSL.}
Unlabeled sets collected automatically at scale are by construction \emph{uncurated}, meaning their contents (features and true labels) are intended to be similar to the target labeled set but not carefully validated.
When the unlabeled set contains images from classes other than the $C$ classes represented in the labeled set, others call this ``open-set'' SSL~\citep{yuMultiTaskCurriculumFramework2020}. 
More formally, if we were to apply labels to the unlabeled set $\mathcal{D}^U$, the set of such labels may include an unknown number of classes beyond the $C$ known classes in the labeled set, and in the worst case may not even include any examples from some (or all) known classes.
Open-set SSL is a special case of uncurated SSL, because a truly uncurated dataset may also differ (usually slightly) in \emph{feature distributions} from the labeled set.
Our CIFAR evaluation focuses on open-set SSL, our Heart2Heart unlabeled set (Sec.~\ref{sec:results-tmed}) is truly uncurated.
% from medical applications.

\citet{oliverRealisticEvaluationDeep2018}~designed seminal experiments using CIFAR-10 images that purposefully build an unlabeled set (some animals, some not) that is mismatched in class composition from the labeled set (all animals).
As mismatch increases, many SSL methods (e.g. VAT or Pi-Model) score worse than a labeled-only baseline that ignores the unlabeled set.
Our later experiments confirm this (Fig.~\ref{fig:results-cifar10} left).


Several approaches have tried to remedy this deterioration, striving to be robust to open-set unlabeled data.
Such methods are also called \textit{safe SSL}~\citep{guoSafeDeepSemiSupervised2020}, because their goal is to perform no worse than labeled-set-only methods.
%UASD~\citep{chenSemiSupervisedLearningClass2020} filters out unlabeled examples which appear to be ``out of distribution'' according to confidence signals derived from its probabilistic classifier. 
%Alternatively, DS3L~\citep{guoSafeDeepSemiSupervised2020} learns a weighting function that can downweight unlabeled examples to prevent poor labeled-set performance.
Tab.~\ref{tab:related_work_comparison} summarizes previous works, with further discussion below.
These methods have been evaluated primarily on artificially mismatched remixes of datasets like CIFAR, and not yet on uncurated medical data.
% from potential SSL application domains like medicine. 

% Broadly, these existing safe SSL approaches do not perform as well as needed (see Fig.~\ref{fig:results-cifar10} far right).
%\todo{ DELETE? shows that the recent non-safe FixMatch dominates 3 ``safe'' SSL methods (UASD, DS3L, and MTCF).
% we remedy in our later medical evaluations.




